\documentclass[11pt]{article}
\usepackage[margin=1in]{geometry}
\usepackage{amsmath,amssymb,amsthm,mathtools}
\usepackage{enumitem}
\usepackage{booktabs}

\newtheorem{theorem}{Theorem}
\newtheorem{lemma}{Lemma}
\newtheorem{proposition}{Proposition}
\newtheorem{corollary}{Corollary}
\theoremstyle{definition}
\newtheorem{assumption}{Assumption}
\theoremstyle{remark}
\newtheorem{remark}{Remark}

\newcommand{\E}{\mathbb{E}}
\renewcommand{\Pr}{\mathbb{P}}
\newcommand{\sgn}{\operatorname{sgn}}

\title{Exact Monotonicity and Optimal Sample Size in Best-of-$N$ Selection:\\ Proofs}
\author{}
\date{}

\begin{document}
\maketitle

\section{Setup}

Let $(U_i,S_i)_{i\ge 1}$ be i.i.d.\ pairs, with $U_i$ a bounded true utility and
$S_i$ a verifier score taking values $s_1<s_2<\dots<s_K$. Write
\[
p_j=\Pr(S=s_j),\qquad m_j=\E[U\mid S=s_j],\qquad F_j=p_1+\dots+p_j ,\qquad F_0=0,\ F_K=1 .
\]
Best-of-$N$ returns $I_N\in\arg\max_{i\le N}S_i$ and we study
$V(N):=\E[U_{I_N}]$ for integers $N\ge 1$, with forward difference
$\Delta(N):=V(N+1)-V(N)$.

\begin{assumption}[Full support]\label{as:supp}
$p_j>0$ for all $j$, so $0<F_1<F_2<\dots<F_{K-1}<1$.
\end{assumption}

\begin{assumption}[Score-only tie-breaking]\label{as:tie}
$I_N$ is a function of $(S_1,\dots,S_N)$ and of a randomisation $\xi$ independent of
$(U_i,S_i)_{i\le N}$. (Uniform-at-random and lowest-index tie-breaking both qualify.)
\end{assumption}

Assumption~\ref{as:tie} is what makes the problem tractable and it is not cosmetic:
if ties were broken using $U$ (e.g.\ ``among top-scored answers keep the best one'')
the identity below fails.

\section{Reduction to the maximum score}

\begin{lemma}[Reduction]\label{lem:red}
Let $M_N=\max_{i\le N}S_i$ and $m(s_j)=m_j$. Under Assumption~\ref{as:tie},
$V(N)=\E[m(M_N)]$.
\end{lemma}

\begin{proof}
Let $\mathcal G=\sigma(S_1,\dots,S_N,\xi)$. Because the pairs are independent and $\xi$
is independent of everything, $\E[U_i\mid\mathcal G]=\E[U_i\mid S_i]=m(S_i)$ for each $i$.
Since $\{I_N=i\}\in\mathcal G$,
\[
\E[U_{I_N}\mid\mathcal G]=\sum_{i=1}^N \mathbf 1\{I_N=i\}\,\E[U_i\mid\mathcal G]
=\sum_{i=1}^N \mathbf 1\{I_N=i\}\,m(S_i)=m(S_{I_N})=m(M_N).
\]
Take expectations.
\end{proof}

\begin{lemma}[Exact formula]\label{lem:formula}
Let $d_j:=m_{j+1}-m_j$ for $j=1,\dots,K-1$. Then
\begin{align}
V(N)&=\sum_{j=1}^{K}m_j\bigl(F_j^{N}-F_{j-1}^{N}\bigr)
     = m_K-\sum_{j=1}^{K-1}d_j\,F_j^{N}, \label{eq:V}\\
\Delta(N)&=\sum_{j=1}^{K-1}d_j\,(1-F_j)\,F_j^{N}. \label{eq:D}
\end{align}
\end{lemma}

\begin{proof}
$\Pr(M_N\le s_j)=F_j^N$ by independence, so $\Pr(M_N=s_j)=F_j^N-F_{j-1}^N$, and the first
equality is Lemma~\ref{lem:red}. Summation by parts with $F_0^N=0$, $F_K^N=1$ gives
$\sum_{j=1}^K m_j(F_j^N-F_{j-1}^N)=m_K F_K^N-\sum_{j=1}^{K-1}(m_{j+1}-m_j)F_j^N$.
Finally $F_j^{N}-F_j^{N+1}=(1-F_j)F_j^{N}$ gives \eqref{eq:D}.
\end{proof}

For $K=3$, writing $a=F_2=p_1+p_2$ and $b=F_1=p_1$, \eqref{eq:V}--\eqref{eq:D} read
\begin{equation}\label{eq:three}
V(N)=m_3+(m_2-m_3)a^N+(m_1-m_2)b^N,\qquad
\Delta(N)=(m_2-m_1)(1-b)\,b^N-(m_2-m_3)(1-a)\,a^N ,
\end{equation}
which are exactly the two expressions in the project description.

\section{When is Best-of-$N$ monotone for every score distribution?}

\begin{theorem}[Aligned verifier]\label{thm:aligned}
\begin{enumerate}[label=(\roman*)]
\item If $m_1\le m_2\le\dots\le m_K$, then $V$ is nondecreasing in $N$, for every $p$.
It is strictly increasing under Assumption~\ref{as:supp} unless all $m_j$ are equal.
\item Conversely, if $m_{j+1}<m_j$ for some $j$, there is a full-support $p$ for which
$V(2)<V(1)$. Hence $V$ is nondecreasing for every score distribution \emph{iff} $m$ is
nondecreasing.
\end{enumerate}
\end{theorem}

\begin{proof}
(i) Every term of \eqref{eq:D} is $\ge 0$, and under Assumption~\ref{as:supp} the term $j$
is $>0$ whenever $d_j>0$. (Equivalently: $M_{N+1}\ge M_N$ pointwise, so
$m(M_{N+1})\ge m(M_N)$ pointwise; this version needs no discreteness and holds for
arbitrary score distributions.)

(ii) For $\varepsilon\in(0,1)$ put mass $(1-\varepsilon)/2$ on each of $s_j,s_{j+1}$ and
spread $\varepsilon$ evenly over the remaining levels. $\Delta(1)$ is continuous in $p$, and
at $\varepsilon=0$ we have $F_i=0$ for $i<j$, $F_j=\tfrac12$, $F_i=1$ for $i>j$, so
\eqref{eq:D} gives $\Delta(1)=d_j/4<0$. Hence $\Delta(1)<0$ for small $\varepsilon>0$.
\end{proof}

\section{Complete characterization for three levels}

Throughout this section $K=3$, Assumption~\ref{as:supp} holds (so $0<b<a<1$), and
\[
\alpha:=m_2-m_1,\qquad \beta:=m_3-m_2,\qquad r:=b/a\in(0,1).
\]
By \eqref{eq:three},
\begin{equation}\label{eq:g}
\Delta(N)=a^{N}g(N),\qquad g(x):=\alpha(1-b)\,r^{x}+\beta(1-a),\quad x\in\mathbb R .
\end{equation}

\begin{lemma}[Single crossing]\label{lem:sc}
$\sgn\Delta(N)=\sgn g(N)$. If $\alpha>0$, $g$ is strictly decreasing; if $\alpha<0$, strictly
increasing; if $\alpha=0$, constant. In all cases $g(x)\to\beta(1-a)$ as $x\to\infty$.
Consequently $N\mapsto\sgn\Delta(N)$ changes at most once, and if $\beta\neq 0$ then
$\sgn\Delta(N)=\sgn\beta$ for all large $N$.
\end{lemma}

\begin{proof}
$a^N>0$, $r^x$ is strictly decreasing with limit $0$, and $(1-b)>0$, $(1-a)>0$.
\end{proof}

\begin{theorem}[Two-point monotonicity test]\label{thm:two}
Let $\Delta(1)=V(2)-V(1)=\alpha\,p_1(p_2+p_3)+\beta\,(p_1+p_2)p_3$. Then
\begin{enumerate}[label=(\alph*)]
\item $V$ is nondecreasing on $\{1,2,\dots\}$ $\iff$ $\beta\ge 0$ and $\Delta(1)\ge 0$;
\item $V$ is nonincreasing on $\{1,2,\dots\}$ $\iff$ $\beta\le 0$ and $\Delta(1)\le 0$;
\item otherwise $V$ is not monotone, and exactly one of the following holds:
  \begin{itemize}
  \item $\beta<0<\Delta(1)$: $V$ strictly increases, then strictly decreases
        (interior maximum; ``more sampling eventually hurts'');
  \item $\beta>0>\Delta(1)$: $V$ strictly decreases, then strictly increases (interior minimum).
  \end{itemize}
\end{enumerate}
In words: Best-of-$N$ is monotone iff the sign of the first step $V(2)-V(1)$ agrees with the
sign of the top increment $m_3-m_2$.
\end{theorem}

\begin{proof}
(a) $V$ is nondecreasing iff $g(N)\ge 0$ for all integers $N\ge1$.
If $\alpha>0$: $g$ decreases to $\beta(1-a)$, which it never attains, so
$\inf_{N\ge 1}g(N)=\beta(1-a)$ and $g(N)\ge0\ \forall N$ iff $\beta\ge0$; in that case
$\Delta(1)>0$ automatically.
If $\alpha<0$: $g$ is increasing, so $g(N)\ge0\ \forall N\ge1$ iff $g(1)\ge0$, i.e.\
$\Delta(1)\ge0$; and then $\beta(1-a)=\lim g>g(1)\ge0$ gives $\beta>0$.
If $\alpha=0$: $g\equiv\beta(1-a)$ and $\Delta(1)=\beta(1-a)a$.
In each case the condition is ``$\beta\ge0$ and $\Delta(1)\ge0$''.
(b) Apply (a) to $-U$.
(c) If neither (a) nor (b) holds then either $\beta<0<\Delta(1)$ or $\beta>0>\Delta(1)$
(the remaining sign patterns, e.g.\ $\beta=0$, fall under (a) or (b) by the case analysis
above). In the first case $g(1)>0$ and $g(\infty)<0$ force $\alpha>0$, so $g$ is strictly
decreasing with a single sign change from $+$ to $-$; the second case is symmetric.
\end{proof}

\begin{theorem}[Optimal $N$ in the one-tail-crossover regime]\label{thm:opt}
Suppose $m_1<m_2$ and $m_3<m_2$ (i.e.\ $\alpha>0>\beta$) and define
\begin{equation}\label{eq:Nstar}
N^\star=\frac{\ln\!\Bigl(\dfrac{(m_2-m_1)(1-b)}{(m_2-m_3)(1-a)}\Bigr)}{\ln(a/b)}
=\frac{\ln\!\Bigl(\dfrac{(m_2-m_1)(p_2+p_3)}{(m_2-m_3)\,p_3}\Bigr)}{\ln\!\bigl((p_1+p_2)/p_1\bigr)} .
\end{equation}
Then $\Delta(N)>0$ for $N<N^\star$, $\Delta(N)=0$ for $N=N^\star$, and $\Delta(N)<0$ for
$N>N^\star$. Consequently:
\begin{enumerate}[label=(\roman*)]
\item $V$ is strictly increasing on $\{1,\dots,N_{\mathrm{opt}}\}$ and nonincreasing afterwards
(strictly decreasing from $N_{\mathrm{opt}}+1$ on), where
\[
N_{\mathrm{opt}}=\max\bigl\{1,\ \lfloor N^\star\rfloor+1\bigr\}
\]
is a maximizer of $V$ over $N\ge1$. It is the unique maximizer unless $N^\star$ is an
integer $\ge1$, in which case the maximizers are exactly $N^\star$ and $N^\star+1$.
\item The optimum is interior ($N_{\mathrm{opt}}\ge2$) iff $N^\star\ge 1$ iff $\Delta(1)\ge 0$, i.e.
\[
(m_2-m_1)\,p_1(p_2+p_3)\ \ge\ (m_2-m_3)\,(p_1+p_2)\,p_3 ,
\]
with $V(2)>V(1)$ iff the inequality is strict; otherwise $N_{\mathrm{opt}}=1$ and sampling
more is harmful from the first extra sample.
\item $V(N)\to m_3<V(N_{\mathrm{opt}})$ as $N\to\infty$.
\end{enumerate}
\end{theorem}

\begin{proof}
By \eqref{eq:g}, with $c:=\dfrac{(m_2-m_3)(1-a)}{(m_2-m_1)(1-b)}>0$,
\[
\Delta(N)>0\iff \alpha(1-b)r^{N}>-\beta(1-a)\iff r^{N}>c\iff N\ln r>\ln c
\iff N<\frac{\ln c}{\ln r}=\frac{\ln(1/c)}{\ln(a/b)}=N^\star,
\]
where the last equivalence divides by $\ln r<0$. The same chain with ``$=$'' and ``$<$''
gives the other two cases.

(i) If $N^\star<1$, $\Delta(N)<0$ for all $N\ge 1$ and the maximizer is $N=1=\max\{1,\lfloor N^\star\rfloor+1\}$.
If $N^\star\ge1$, the integers with $\Delta(N)>0$ are exactly $N<N^\star$.
When $N^\star\notin\mathbb Z$ these are $N\le\lfloor N^\star\rfloor$, so $V(1)<\dots<V(\lfloor N^\star\rfloor+1)$
and $V$ strictly decreases afterwards. When $N^\star\in\mathbb Z$, $V$ strictly increases up to
$V(N^\star)$, $V(N^\star+1)=V(N^\star)$, and $V$ strictly decreases afterwards.

(ii) $N^\star\ge1\iff r\ge c\iff g(1)\ge 0\iff\Delta(1)\ge0$; substitute $1-a=p_3$, $1-b=p_2+p_3$.

(iii) $a^N,b^N\to0$ in \eqref{eq:three}. For the strict inequality, pick $N>N^\star$, $N\ge N_{\mathrm{opt}}$:
$V(N)-m_3=\sum_{n\ge N}(-\Delta(n))>0$.
\end{proof}

\begin{remark}[Reading \eqref{eq:Nstar}]
$N^\star$ grows when the benefit of leaving the bottom level, $(m_2-m_1)$, is large relative
to the over-optimization penalty $(m_2-m_3)$; when the misleading top level is rare
($p_3\to0$ gives $N^\star\to\infty$); and when $p_1$ is close to $p_1+p_2$ so that the
``good'' term $b^N$ decays almost as slowly as the ``bad'' term $a^N$.
\end{remark}

\section{Generalization to $K$ levels}

For $K\ge4$, $\Delta$ is a sum of more than two exponentials and need not cross zero only
once. The right generalization bounds the number of turning points of $V$ by the number of
turning points of the calibration curve $j\mapsto m_j$.

\begin{lemma}[Descartes' rule for exponential sums]\label{lem:desc}
Let $\lambda_1<\dots<\lambda_n$ and let $c_1,\dots,c_n\in\mathbb R$ be not all zero. Then
$f(x)=\sum_{j=1}^n c_j e^{\lambda_j x}$ has at most $Z$ real zeros counted with multiplicity,
where $Z$ is the number of sign changes in $(c_1,\dots,c_n)$ after deleting zeros.
\end{lemma}

\begin{proof}
Induct on $Z$. If $Z=0$, all nonzero $c_j$ share a sign and $f$ has no zero. If $Z\ge1$,
delete the zero coefficients, pick consecutive indices $k,k+1$ with $c_kc_{k+1}<0$ and
$\mu\in(\lambda_k,\lambda_{k+1})$. Then $h(x)=e^{-\mu x}f(x)$ has the same zeros with the same
multiplicities as $f$, and
$h'(x)=\sum_j c_j(\lambda_j-\mu)e^{(\lambda_j-\mu)x}$. The factors $(\lambda_j-\mu)$ are negative
for $j\le k$ and positive for $j\ge k+1$, so the coefficient sequence of $h'$ has exactly
$Z-1$ sign changes. By induction $h'$ has at most $Z-1$ zeros, and by Rolle's theorem
(with multiplicity) $h$ has at most $Z$.
\end{proof}

\begin{theorem}[Variation diminishing]\label{thm:K}
Under Assumptions \ref{as:supp}--\ref{as:tie}, suppose the $m_j$ are not all equal, and let
$Z$ be the number of sign changes in the increment sequence $(d_1,\dots,d_{K-1})$ (zeros deleted). Then:
\begin{enumerate}[label=(\roman*)]
\item the sequence $\Delta(1),\Delta(2),\dots$ changes sign at most $Z$ times, i.e.\
$V$ has at most $Z$ turning points;
\item for all large $N$, $\sgn\Delta(N)=\sgn d_{j^\dagger}$ where $j^\dagger$ is the largest
index with $d_j\ne0$: \emph{the top of the calibration curve decides the tail};
\item \emph{(single-peaked calibration $\Rightarrow$ single-peaked BoN)} if
$m_1\le\dots\le m_k\ge\dots\ge m_K$ with $m_k>m_K$, then $V$ is nondecreasing on
$\{1,\dots,N_{\mathrm{opt}}\}$ and nonincreasing afterwards, where
$N_{\mathrm{opt}}=\min\{N\ge1:\Delta(N)\le0\}<\infty$ maximizes $V$.
\end{enumerate}
\end{theorem}

\begin{proof}
By \eqref{eq:D}, $\Delta(N)=f(N)$ with $f(x)=\sum_{j=1}^{K-1}c_je^{\lambda_jx}$,
$c_j=d_j(1-F_j)$, $\lambda_j=\ln F_j$. Assumption~\ref{as:supp} gives
$\lambda_1<\dots<\lambda_{K-1}<0$ and $\sgn c_j=\sgn d_j$.

(i) By Lemma~\ref{lem:desc}, $f$ has at most $Z$ real zeros. Each sign change of the integer
sequence $\Delta(N)$ forces, by the intermediate value theorem, a sign change of $f$ on a
distinct interval, hence a distinct zero of odd multiplicity.

(ii) $f(x)e^{-\lambda_{j^\dagger}x}\to c_{j^\dagger}$ as $x\to\infty$, since all other
surviving exponents are strictly smaller.

(iii) Here the nonzero increments are of the form $(+,\dots,+,-,\dots,-)$ with at least one
``$-$'', so $Z\le1$ and by (ii) $f<0$ eventually. If $f$ has no zero, $f<0$ everywhere and
$N_{\mathrm{opt}}=1$. Otherwise $f$ has exactly one zero $x_0$, which is simple (a zero of
multiplicity $\ge 2$ would exceed the bound $Z\le 1$), so $f>0$ on $(-\infty,x_0)$ and $f<0$ on $(x_0,\infty)$.
Either way $\Delta(N)>0$ for $N<N_{\mathrm{opt}}$ and $\Delta(N)\le0$ for $N\ge N_{\mathrm{opt}}$.
\end{proof}

Theorem~\ref{thm:K} contains the earlier results: $Z=0$ with positive increments is
Theorem~\ref{thm:aligned}(i); for $K=3$ it reproduces Lemma~\ref{lem:sc}, and
Theorem~\ref{thm:opt} is the case where the unique zero $x_0=N^\star$ is available in closed form.
For $K\ge4$ there is in general no closed form, but (iii) says a linear scan for the first
$N$ with $\Delta(N)\le0$ (or bisection on $f$) finds the global optimum.

\begin{remark}[Continuous scores]
Lemma~\ref{lem:red} holds for any score distribution. If $S$ has continuous cdf $F$ and
$m$ has bounded variation, \eqref{eq:D} becomes
$\Delta(N)=\int F(s)^N(1-F(s))\,dm(s)$, and the kernel $F(s)^N=e^{N\ln F(s)}$ is totally
positive, so the same variation-diminishing conclusion follows from Karlin's theorem
(S.~Karlin, \emph{Total Positivity}, 1968). We do not need this here: the finite pools in the
experiment below are covered by the finite-$K$ formula (Proposition~\ref{prop:rank}).
\end{remark}

\section{Connection to the LLM experiment}

The experiment draws, for each question $q$, a pool of $M$ scored answers and averages
Best-of-$N$ accuracy over questions. Questions are heterogeneous, and an average of
single-peaked curves need not be single-peaked, so one might worry that the theory does not
transfer. The \emph{formula} does, because ranks are quantiles.

\begin{proposition}[Rank calibration]\label{prop:rank}
For each question let $c_{q,j}\in\{0,1\}$ be the correctness of its $j$-th lowest-scored answer
(scores distinct), and $\bar m_j=\frac1Q\sum_q c_{q,j}$ the \emph{rank-calibration curve}.
\begin{enumerate}[label=(\roman*)]
\item If Best-of-$N$ draws i.i.d.\ (with replacement) from each question's pool, the
question-averaged accuracy is exactly \eqref{eq:V} with $K=M$, $F_j=j/M$, $m_j=\bar m_j$.
So Theorems \ref{thm:aligned} and \ref{thm:K} apply verbatim to $\bar m$.
\item If it draws without replacement (what \texttt{run\_experiment.py} does),
\[
V(N)=\sum_{j=N}^{M}\bar m_j\,\binom{j-1}{N-1}\Big/\binom{M}{N},
\]
and $V$ is nondecreasing whenever $\bar m$ is nondecreasing.
\end{enumerate}
\end{proposition}

\begin{proof}
(i) For a single question, the pool's empirical distribution is a $K=M$ level model with
$p_j=1/M$ and $m_j=c_{q,j}$, so $V_q(N)=\sum_j c_{q,j}\bigl((j/M)^N-((j-1)/M)^N\bigr)$. The
weights do not depend on $q$ and $V_q$ is linear in $c_{q,\cdot}$; average over $q$.
(ii) The top-scored element of a uniform $N$-subset has rank $j$ iff $j$ is chosen together
with $N-1$ of the $j-1$ lower ranks. Monotonicity: a uniform $(N+1)$-subset is a uniform
$N$-subset plus one more element, so its maximum rank is pointwise no smaller.
\end{proof}

\begin{remark}[What this does and does not test]\label{rem:identity}
Proposition~\ref{prop:rank} is an identity: it reproduces \emph{any} observed curve from that
curve's own rank calibration, so agreement with it is not evidence for the theory. Nor does
Theorem~\ref{thm:K} restrict the raw curves much: at $M=256$ the empirical $\bar m$ has more than
$150$ turning points, almost all sampling noise. What the data can test is whether a real
scorer has an interior optimum, and whether a coarse three-level model predicts it.
Tied scores are common (the Skywork scorers run in half precision); tied answers are given their
mean correctness, which is uniform tie-breaking (Assumption~\ref{as:tie}).
\end{remark}

\subsection{Data}

Unless noted, answers come from Qwen2.5-Math-1.5B-Instruct ($M=256$ per question, temperature $0.8$, top-$p$
$0.95$) and are scored by a reward model: Skywork-Reward-V2-Qwen3-0.6B (``Skywork-0.6B'', the weak
scorer), its 1.7B version (``Skywork-1.7B''), or OpenAssistant DeBERTa-v3-large (``DeBERTa''), a general-chat model not
trained on mathematics. The weak generator is Qwen2.5-0.5B-Instruct. The held-out MATH set is a random $3{,}000$ of the
$4{,}500$ MATH test questions that are not in MATH-500. All curves below are exact, and all intervals are $95\%$
bootstrap intervals over questions; the bootstrap does not account for choosing a peak.

\begin{center}
\small
\begin{tabular}{lrrccccr}
\toprule
run & $Q$ & $M$ & $V(1)$ & $V(16)$ & $V(64)$ & $V(M)$ & $\arg\max_N$ \\
\midrule
GSM8K, Skywork-1.7B                   & 100  & 64  & .859 & .922 & .950 & .950 & 64 \\
GSM8K, Skywork-0.6B                   & 100  & 64  & .859 & .908 & .920 & .920 & 64 \\
GSM8K, Skywork-0.6B, weak generator   & 100  & 64  & .439 & .670 & .740 & .740 & 64 \\
GSM8K, Skywork-0.6B, first 100        & 100  & 256 & .852 & .911 & .928 & .920 & 95 \\
GSM8K, Skywork-0.6B, all questions    & 1319 & 256 & .853 & .899 & .903 & .909 & 256 \\
MATH-500, Skywork-0.6B                & 500  & 256 & .749 & .805 & .810 & .798 & 70 \\
MATH held-out, Skywork-0.6B           & 3000 & 256 & .758 & .812 & .811 & .808 & 24 \\
MATH-500, DeBERTa                     & 500  & 256 & .749 & .738 & .727 & .702 & 2 \\
\bottomrule
\end{tabular}
\end{center}

\subsection{Three shapes, and no confirmed interior optimum}

\begin{itemize}[leftmargin=*]
\item \textbf{Rising.} On all $1{,}319$ GSM8K questions $V$ rises throughout, from $.853$ to $.909$,
though only slowly after $N=16$: $V(256)-V(16)=.0100$, interval $[.0004,.0194]$.
\item \textbf{Rise, then plateau.} On held-out MATH, Skywork-0.6B gains about five points by $N=16$ and then is
flat. Its best $N$ is $24$, but $V(24)-V(256)=.0041$ with interval $[-.003,.011]$.
\item \textbf{Falling.} On MATH-500, DeBERTa gives $V(1)-V(256)=.0473$, interval $[.0232,.0719]$. The small
rise to $N=2$ ($.0020$, interval $[-.0010,.0050]$) is indistinguishable from zero, so $N=1$ and $N=2$ cannot
be separated, and the curve is effectively nonincreasing.
\end{itemize}

\paragraph{Both apparent peaks failed to replicate.}
Two runs showed an interior peak on a small sample. Neither survived more questions.

\begin{center}
\small
\small
\begin{tabular}{lll}
\toprule
peak seen in & decline to $N=256$ & same comparison, new questions \\
\midrule
GSM8K, first 100 ($N=94$) & $.0083$, $[-.0104,.0308]$ & $-.0053$, $[-.0119,.0011]$ (1{,}219 questions) \\
MATH-500 ($N=64$) & $.0119$, $[-.0017,.0255]$ & $\phantom{-}.0028$, $[-.0030,.0086]$ (3{,}000 questions) \\
\bottomrule
\end{tabular}
\end{center}

\noindent Neither original decline was significant either. The second row is the only confirmatory test in this note:
$N=64$ and the sample size were fixed, and committed to the repository, before the held-out data
were generated. ($N=64$ was itself picked from the MATH-500 curve.) The design had about $94\%$ power for a true decline of one point. The held-out interval excludes the MATH-500 point
estimate of $.0119$, but not a decline of up to about $.009$. The first row was not planned in advance.
The practical lesson is that the argmax of a noisy curve on a few hundred questions readily produces a false peak.

\subsection{Does a three-level model predict the curve?}

We fit three-level models with within-question score bins of sizes (bottom, middle, top), using $\bar m$ from the same data, and applied
Theorems \ref{thm:two}--\ref{thm:opt}. The $13$ binnings are tertiles, plus bottom $\in\{1/2,3/4,7/8,15/16\}$
of the answers crossed with a top level of $1$, $2$ or $4$ answers. They were chosen after looking at the data. The grid
includes all four $K=3$ binnings we had tried on MATH-500 with a top level of one or two answers, and tertiles; equal bins with $K=2$ and $K=4$, and other
splits on the first $100$ GSM8K questions, were also tried.

\begin{center}
\small
\begin{tabular}{lccccc}
\toprule
& \multicolumn{3}{c}{predicted shape (13 binnings)} & $N_{\mathrm{opt}}$ & $V(256)$ \\
run ($M=256$) & increasing & peak & decreasing & if peaked & fit vs.\ data \\
\midrule
GSM8K, all questions   & 13 & 0 & 0 & --      & .891--.907 vs.\ .909 \\
MATH-500, Skywork      & 10 & 3 & 0 & 28--76  & .794--.814 vs.\ .798 \\
MATH held-out, Skywork & 5  & 8 & 0 & 25--86  & .804--.811 vs.\ .808 \\
MATH-500, DeBERTa      & 0  & 4 & 9 & 2--4    & .715--.749 vs.\ .702 \\
\bottomrule
\end{tabular}
\end{center}

\noindent The fits get clear shapes right (GSM8K rises, DeBERTa essentially falls from the start). They do not
locate $N_{\mathrm{opt}}$: where the top of the calibration curve is nearly flat, both the presence of a peak and
its position depend on the binning. The fitted top-level drop $m_2-m_3$ in the held-out peaked fits is at most
$.007$. The level of the curve is also off: every binning overestimates DeBERTa's $V(256)$ by at least $.013$,
because the drop sits in the very top ranks (the single top answer is right for $70.2\%$ of questions,
against $74.9\%$ for a random one) and coarse bins average it away. Tertiles, the natural default, miss the
DeBERTa drop entirely (predicting $V(256)=.749$). These are in-sample fits; no held-out prediction of $N_{\mathrm{opt}}$ was attempted.

\subsection{What was tried, and limitations}

\begin{itemize}[leftmargin=*]
\item \textbf{Exploration.} Besides the runs above, we blended the Skywork and DeBERTa scores on
MATH-500 at ten weights, looking for an interior peak; none showed both a significant rise and a significant
decline, and the decline did not vary steadily with the weight. A cross-fitting analysis was discarded because resampled copies of a question could fall in both halves. The DeBERTa comparison was also computed from $N=16$, $64$ and $94$, all with intervals
above zero; no multiplicity correction was applied.
\item \textbf{DeBERTa rests on one dataset.} It has not been replicated. DeBERTa reads at most $512$ tokens, so longer answers were cut from the front,
and part of its weakness may come from truncation.
\item \textbf{Scope.} Three scorers, two generators and two benchmarks, all with $M\le256$: a peak beyond
$N=256$, or for a scorer of intermediate quality, is not excluded.
\end{itemize}

\paragraph{Summary.}
The theory says exactly when a three-level verifier produces a rising, falling or single-peaked
curve. In the data, a scorer with useful, nearly flat-topped calibration rises to a plateau, and a scorer with no useful
calibration and a bad top falls. We found no confirmed instance of the one-tail-crossover regime in between.
That regime remains a theoretical possibility these experiments neither demonstrate nor rule out.

\end{document}
